\documentclass[journal]{IEEEtran}

\usepackage[T1]{fontenc}
\usepackage{amsmath,amsfonts}
\usepackage{newtxtext}
\usepackage{newtxmath}

\usepackage{algorithmic}

\usepackage{cite}
\usepackage{graphicx}
\usepackage{booktabs}
\usepackage{tabularx}
\usepackage{geometry}

\usepackage[hidelinks]{hyperref}

\begin{document}

\title{A Methods Plan to Test a Wearable Anxiety Monitor}

\author{Andy Babcock \\
    \textit{Steve Sanghi College of Engineering, Northern Arizona University} \\
    Flagstaff, AZ, USA \\
    aab726@nau.edu
}

\maketitle
\section{Study Methods}
Expanding upon last week's initial device design, much more detail on the specific methods is needed to properly focus this assignment. Previously, the tentative plan was to use a consumer wearable electroencephalogram (EEG) to measure anxiety levels during test-taking. The EEG selected is the NeuroSky MindWave Mobile 2, but the selection of test-taking proved a poor match for the scope of available participants for this study. With the main participants of any study conducted being myself and Cameron Brebner, it was not feasible to create an extended bank of test questions that would provide a difficult challenge and naturally induce anxiety. Instead, adapting the Trier Social Stress Test (TSST) to work between two people is much more achievable~\cite{kirschbaum_pirke_hellhammer_1993}. TSST traditionally was a way to induce anxiety in participants during an anticipation period and a test period though creating, then giving, a presentation of a random topic. However, at the end of the anticipation period (creating the presentation), the notes would be unexpectedly taken away, and the participant forced to give the 5-minute presentation blind. While this testing format works with fewer people, it can be adjusted and scaled to work with only two participants. Each of the two participants will create presentation notes and accompanying slides on something they are knowledgeable in, and the other participant will present on that topic to the first participant. Then, at a random time during the presentation, the notes will be taken away, and the presenter forced to finish the presentation using only what is on the slides themselves.

It is difficult to properly re-create TSST's anticipation period with a small number of participants, and no dedicated judges to add pressure onto the participants while creating presentations. Currently, this step does not induce anxiety aside from both participants knowing that they will need to trade presentations. However, it is difficult to induce anxiety without it being artificial, like playing stressful or unsettling music. However, taking EEG measurements during this phase can instead provide baseline brainwave readings while working in a low-stress environment. The slides would be created on a laptop which could be plugged into a projector or TV, while the notes written down on paper or printed out, making it easy to take them away later.

For the equivalent of TSST's test period, trading slide-decks and presenting blind mimics some of the induced stress from the original test. However, the participant giving the presentation will also be given the presentation notes at the beginning of the 5-minute presentation (5 minutes being chosen as an uncomfortable, but not completely unreasonable amount of time to present; also, because that is what TSST uses). This format loses TSST's surprise factor of the participant losing their notes at the beginning of the presentation, so it is added back in the form of taking away the presenting notes after a random amount of time. This reintroduces an unknown variable, even with prior knowledge, and creates two separate scenarios to measure anxiety under. While presenting with notes, it is stressful but possible thanks to said notes. After taking them away, the presenter is in a similar situation to the presenter from the TSST test, forced to present on something they are not comfortable with while not having notes.

If this study could be scaled up to a larger number of participants, it would be much more practical to use TSST as it was initially created, as the increased number of participants allows for it to run smoothly. However, this structure allows for a middle-ground to induce and measure anxiety with a small number of participants.

\section{Collecting Data}
During both phases, both participants will wear the EEG to record brain activity (for the preparation section, both participants can record brain activity by preparing slides and notes one at a time). The MindWave Mobile 2 records all of the data needed for anxiety analysis, as shown in Table~\ref{tab:eegid_brainwaves}. To record this data, the headset will be paired to a smartphone, and the data synced post-hoc to a laptop to perform data analysis. The chosen phone is a Motorola Moto X4, as it is about the same age as the headset and companion app, eegID~\cite{gsm_moto_x4_2017,neurosky2026eegid}. When attempting to install eegID on modern smartphones, there are significant issues with installing the older mobile app, and with reliably pairing the headset, which uses an older Bluetooth version. After recording brain activity through eegID, the CSV file is synced to an external computer for analysis, likely using KDE Connect to send the file to a paired computer~\cite{kde_connect_2026}.

All of the collected bands contain useful information on various components of anxiety. Notably, the Beta band can be subdivided into multiple categories associated with varying levels of anxiety, from none, slight, mild, moderate, and severe anxiety~\cite{chatterjee_gavas_saha_2023}. Additionally, Theta waves are correlated with relaxation, and can become erratic while under stress~\cite{neurohealth_2019}. Gamma waves can also indicate agitation when the brain is over-aroused and highly anxious. Lastly, Alpha waves indicate a lack of anxiety, so unstable Alpha waves would indicate an anxiety response. There is a multitude of interesting literature on how to interpret specific bands of frequencies, and how to compare them to others, like Wei et al's study on how anxiety shapes theta band activity~\cite{wei_sun_long_2025}. However, adding much beyond a simple 4-level anxiety classification with some extra nuance and context from other brain activity would likely be beyond the scope of this study.

To apply these brainwave categorizations, relatively simple python scripting will suffice. Some of the more nuanced indicators, like disturbances in Alpha waves, may require comparing the variance in stressful and stress-free data, but the main classifications done will be relatively simple. The main hurdle will be accounting for error, from both the headset being cheap and possibly imprecise relative to medical-grade EEGs, and from the readings being taken through skin and bone. Chatterjee's study on training an AI on existing anxiety classifications from brain activity may prove a vital component while analyzing data, as the data may be somewhat unusable in the state it gets collected in~\cite{chatterjee_gavas_saha_2023}.

\begin{table*}[t]
    \centering
    \renewcommand{\arraystretch}{1.4}
    \begin{tabularx}{\textwidth}{@{}llX@{}}
        \toprule
        \textbf{eegID Field} & \textbf{Frequency Range} & \textbf{Description and Meaning} \\
        \midrule
        PoorSignal & N/A & A hardware metric indicating the quality of the sensor contact and signal integrity. \\
        EEG Raw Value & N/A & Unfiltered analog values captured directly from the EEG sensor. \\
        EEG Raw Value Volts & N/A & The raw EEG data scaled and converted into standard voltage units. \\
        Attention Level & N/A & A computed metric (ThinkGear algorithm) measuring the user's focus, concentration, and directed mental effort. \\
        Meditation Level & N/A & A computed metric (ThinkGear algorithm) measuring the user's level of mental calmness and relaxation. \\
        Blink Strength & N/A & A metric that detects eye blinks and quantifies their relative physical intensity. \\
        Delta & 1--3 Hz & Associated with deep, dreamless sleep, trance, and the unconscious mind. Represents a decreased awareness of the physical world. \\
        Theta & 4--7 Hz & Linked to intuition, creativity, daydreaming, internal focus, and spiritual awareness. Reflects the state between wakefulness and sleep. \\
        Alpha Low & 8--9 Hz & Correlates to an inner-awareness of self, mind/body integration, and balance. \\
        Alpha High & 10--12 Hz & Associated with centering, healing, and bridging the conscious to the subconscious. Promotes mental resourcefulness and calmness. \\
        Beta Low & 13--17 Hz & Reflects a state of being relaxed yet focused, integrated, and actively thinking while remaining aware of surroundings. \\
        Beta High & 18--30 Hz & Indicates high alertness and mental activity (e.g., math, planning), and general activation of mind/body functions, but can also reflect agitation. \\
        Gamma Low & 31--40 Hz & Consolidates information from different areas of the brain for simultaneous processing. Associated with alertness, good memory, or potential agitation. \\
        Gamma Mid & 41--50 Hz & Functions similarly to Gamma Low in consolidating rapid, simultaneous processing across brain regions; also indicates high alertness or agitation. \\
        \bottomrule
    \end{tabularx}
    \vspace{5pt}
    \caption{eegID Telemetry Fields and Correlated Brainwave States~\cite{neurosky2026eegid,neurohealth_2019}}
    \label{tab:eegid_brainwaves}
\end{table*}

\bibliographystyle{IEEEtran}
\bibliography{references}

\end{document}
